% 张宇八套卷重新提取勘误（由 build/zy8_errata.py 生成，请勿手改）。
%
% 每份 PDF 只放一个科目；条目按「年份-套次 / 题号：原内容 → 正确内容」给，
% 红色 = 题库原录入，绿色 = 更正后内容。
%
% 依据：对源页 PNG 的原生多模态观察（逐条裁片放大核对），
% 未使用 OCR，未使用 PDF 文本层。
\documentclass[11pt,a4paper]{ctexart}

\usepackage[margin=2.2cm,top=2.6cm]{geometry}
\usepackage{amsmath,amssymb,bm}
\usepackage{tikz}
\usetikzlibrary{patterns,arrows.meta,calc}
\usepgflibrary{plothandlers}
\usepackage{xcolor}
\usepackage{fancyhdr}
\usepackage{hyperref}

\definecolor{errata-red}{RGB}{198,40,40}
\definecolor{errata-green}{RGB}{20,120,60}

\newcommand{\rc}[1]{\textcolor{errata-red}{#1}}
\newcommand{\gc}[1]{\textcolor{errata-green}{#1}}
\newcommand{\Paper}[1]{\par\vspace{9pt}{\bfseries #1}\par\vspace{3pt}}
\newcommand{\pair}[1]{\par\hangindent=2.2em\hangafter=1 #1\par\vspace{1pt}}

\newcommand{\ErrataSubject}{数学三}

% 题库正文里使用的两个项目宏。勘误文档不排版整卷，只把它们的**印刷形态**画出来，
% 这样勘误条目与原卷面部一致（选择题的括号、填空题的横线）。
\newcommand{\ExamSelection}{（\quad）}
\newcommand{\ExamBlank}{\underline{\hspace{3em}}}

\pagestyle{fancy}
\fancyhf{}
\renewcommand{\headrulewidth}{0pt}
\fancyfoot[C]{\small\thepage}

\setlength{\parindent}{0pt}
\linespread{1.12}
\hypersetup{colorlinks=true,linkcolor=errata-green,urlcolor=errata-green}

\begin{document}

\begin{center}
  {\LARGE\bfseries 张宇八套卷题库勘误清单（\ErrataSubject）}
\end{center}
\vspace{2pt}
\begin{center}
  \small 依据源页逐题回图核对，共 163 处经复核，其中 13 处为题库识别错误并已更正。
\end{center}
\vspace{6pt}
\hrule
\vspace{8pt}

\Paper{2025 张宇八套卷（数学三）第 3 套}
\pair{2：\rc{（整段重录）} $\longrightarrow$ \gc{设常数$m>0,n>0$，则$\int_0^n\sqrt{x}\left[\dfrac{m}{x}\right]\,\mathrm{d}x$（$[\cdot]$是取整符号）的敛散性\ExamSelection}}
\Paper{2025 张宇八套卷（数学三）第 4 套}
\pair{21：图形重绘（源页 zhangyu8-2025-math3/pages/p018.png）。源页 x5 边的箭头在靠 (100) 交汇点一端、指向右上，即 x5 自下方节点流向右侧节点；题库绘成 (R)--(B)（自右节点流向下方节点），流向画反。其余 a、x1、x2、x3、x4、40、100、b 各边与源页一致}
\Paper{2025 张宇八套卷（数学三）第 6 套}
\pair{7（A）：\rc{（整段重录）} $\longrightarrow$ \gc{$r(A\quad AB)=r(A)$。}}
\pair{7（B）：\rc{（整段重录）} $\longrightarrow$ \gc{$r(AB^{\mathrm T}\quad AB^{\mathrm T}B)=r(AB^{\mathrm T})$。}}
\Paper{2025 张宇八套卷（数学三）第 8 套}
\pair{19：\rc{（整段重录）} $\longrightarrow$ \gc{求 $\displaystyle \sum_{n=1}^{\infty}\dfrac{\left[2+(-1)^n\right]^n}{n\cdot 6^n}$ 的值。}}
\Paper{2026 张宇八套卷（数学三）第 3 套}
\pair{5（A）：\rc{（整段重录）} $\longrightarrow$ \gc{$AP_1^9P_2^{\mathrm T}=B$}}
\pair{5（B）：\rc{（整段重录）} $\longrightarrow$ \gc{$AP_2^{\mathrm T}P_1^9=B$}}
\pair{5（C）：\rc{（整段重录）} $\longrightarrow$ \gc{$P_1^9P_2^{\mathrm T}A=B$}}
\pair{5（D）：\rc{（整段重录）} $\longrightarrow$ \gc{$P_2^{\mathrm T}P_1^9A=B$}}
\Paper{2026 张宇八套卷（数学三）第 4 套}
\pair{3（A）：\rc{（整段重录）} $\longrightarrow$ \gc{$\displaystyle\int_0^1\ln\left[1+\dfrac{f(x)}{n}\right]\mathrm{d}x$。}}
\pair{5（C）：\rc{（整段重录）} $\longrightarrow$ \gc{$\begin{pmatrix}1&0&0\\2&1&0\\3&\dfrac43&1\end{pmatrix}\begin{pmatrix}1&0&0\\0&3&0\\0&0&-\dfrac23\end{pmatrix}\begin{pmatrix}1&-1&2\\0&1&-\dfrac13\\0&0&1\end{pmatrix}$。}}
\pair{5（D）：\rc{（整段重录）} $\longrightarrow$ \gc{$\begin{pmatrix}1&0&0\\2&-1&0\\3&\dfrac43&1\end{pmatrix}\begin{pmatrix}1&0&0\\0&-3&0\\0&0&-\dfrac23\end{pmatrix}\begin{pmatrix}1&-1&2\\0&1&-\dfrac13\\0&0&1\end{pmatrix}$。}}
\pair{15：\rc{（整段重录）} $\longrightarrow$ \gc{设$\boldsymbol{\alpha}_1=\begin{pmatrix}4\\1\\2\end{pmatrix}$，$\boldsymbol{\alpha}_2=\begin{pmatrix}1\\1\\-1\end{pmatrix}$，$\boldsymbol{\alpha}_3=\begin{pmatrix}-2\\0\\4\end{pmatrix}$，$\boldsymbol{\alpha}_4=\begin{pmatrix}7\\2\\-3\end{pmatrix}$，$\boldsymbol{A}=(\boldsymbol{\alpha}_1,\boldsymbol{\alpha}_2,\boldsymbol{\alpha}_3,\boldsymbol{\alpha}_4)$，则$\boldsymbol{A}\boldsymbol{x}=\boldsymbol{\alpha}_2+3\boldsymbol{\alpha}_4$的通解为\ExamBlank。}}

\end{document}
